\section{Grundlagen}
\subsection{Smart Grids: physische Ebene sowie Informations- und Kommunikationstechnik}
Moderne Energiesysteme befinden sich in einem tiefgreifenden Wandel, der durch den verstärkten Einsatz erneuerbarer Energien, dezentrale Erzeugungsstrukturen und eine zunehmende Volatilität im Netzbetrieb geprägt ist. Um diesen Herausforderungen zu begegnen, entwickeln sich klassische Stromnetze zu sogenannten Smart Grids, die physische Infrastruktur mit digitalen Technologien zur Überwachung, Steuerung und Effizienzsteigerung verbinden (Gungor et al., 2013). Die physische Ebene umfasst dabei Erzeugungsanlagen, Umspannwerke, Leitungen und Lasten, während die digitale Ebene aus Messsystemen, Kommunikationsnetzen, Regelalgorithmen und Datenplattformen besteht (Fang et al., 2012).

Eine zentrale Rolle spielt die Advanced Metering Infrastructure (AMI), die durch intelligente Zähler eine Echtzeiterfassung des Energieverbrauchs ermöglicht. Ergänzt wird sie durch SCADA-Systeme, Sensorik sowie standardisierte Kommunikationsprotokolle wie IEC 61850 (NIST, 2014). Diese Vernetzung verbessert Überwachung und Steuerung des Netzes, erhöht jedoch zugleich dessen Verwundbarkeit. Störungen, Softwarefehler oder gezielte Angriffe können sich aufgrund der engen Kopplung von digitaler und physischer Ebene schnell ausbreiten (Buldyrev et al., 2010). Für die Analyse der Widerstandsfähigkeit ist es daher entscheidend, beide Ebenen gemeinsam zu betrachten, da nur ihr Zusammenspiel das Gesamtverhalten des Systems erklärt.

 
\subsection{Komplexe Netzwerke: Knoten, Kanten, typische Topologien}
Komplexe Netze zeigen vereinfacht, wie echte Systeme funktionieren – dabei werden Bausteine mit ihren Verbindungen durch Grafiken dargestellt. Eine solche Struktur enthält Punkte, sogenannte Knoten, verbunden über Linien, genannt Kanten, welche Austauschprozesse abbilden (Newman, 2018). In intelligenten Stromnetzen stehen diese Knoten meist für technische Komponenten wie Umspannwerke, Transformatorstationen, Steuerzentralen oder lokale Energiequellen. Die verbindenden Elemente hingegen symbolisieren entweder reale Leitungen. Entscheidend dafür, wie gut das Gesamtsystem Störungen aushält, ist nicht nur die Menge der beteiligten Einheiten, sondern auch deren Vernetzungsart.
\subsubsection{Knoten}
Ein Knoten bildet die Basis eines Netzwerks. Abhängig vom Kontext handelt es sich um ein Gerät, eine Stelle in einer Struktur oder einen gedanklichen Punkt. Die Zahl der Linien, die von ihm wegführen, heißt Grad des Knotens. Wichtig ist dieser Wert, weil er zeigt, wie bedeutend ein Knoten ist. Besonders viele Verbindungen haben oft wichtige Knoten im Stromnetz, das macht sie empfindlich.
\subsubsection{Kanten} 
Kanten zeigen, wie Knoten miteinander verbunden sind. In technischen Systemen, etwa bei Smart Grids, dienen sie dem Transport von Energie oder Daten. Zu ihren zentralen Merkmalen gehören sogenannte gewichtete Kanten. Dieses Gewicht steht im Stromnetz beispielsweise für Parameter wie Leitungsbelastbarkeit oder Widerstand. Darüber hinaus spielt die Richtung (gerichtet/ungerichtet) eine Rolle. Bei rein struktureller Betrachtung gelten Stromnetze als ungerichtet; sobald Leistungsflüsse ins Gewicht fallen, werden sie dagegen als gerichtet angesehen.
Knoten mit Kanten prägen die Struktur eines Netzwerks, was sein Verhalten bei Belastung beeinflusst, ebenso wie die Reaktion auf Störungen.
\subsubsection{Zentralitätsmaße}
Zentralitätsmaße quantifizieren die Bedeutung eines Knotens. Beispiele dafür sind die folgenden drei Maße.
	Betweenness Centrality: Sie misst, wie häufig ein Knoten auf den kürzesten Pfaden zwischen anderen Knoten liegt. Das ist wichtig für die Flusskontrolle. Sie wird durch folgenden Ausdruck gemessen:
\[
g(v) = \sum_{\substack{s \neq v \neq t}} \frac{\sigma_{st}(v)}{\sigma_{st}}
\]

$\sigma_{st}$ ist dabei die die Gesamtzahl der kürzesten Wege vom Knoten s zum Knoten t und $\sigma_{st}(v)$ die Anzahl derjenigen dieser Wege, die durch den Knoten v verlaufen (wobei v kein Endpunkt des Weges ist).
Closeness Centrality: Sie misst, wie nah ein Knoten durchschnittlich zu allen anderen liegt. Das erfolgt folgendermaßen: 
\[
C(x) = \frac{N - 1}{\sum_{y} d(y, x)}
\]

$N$ ist in der Abbildung dabei die Anzahl der Knoten im Graphen. Die Normalisierung der Closeness (Nähe-Zentralität) erleichtert den Vergleich von Knoten in Graphen unterschiedlicher Größe. Bei sehr großen Graphen wird das „minus eins“ in der Normalisierung vernachlässigbar und daher häufig weggelassen.
Degree Centrality: Die Degree Centrality bezeichnet die Bedeutung eines Knotens in einem Netzwerk anhand der Anzahl seiner direkten Verbindungen. Sie entspricht dem Knotengrad und misst, wie viele Kanten von einem Knoten ausgehen. Ein Knoten mit hoher Degree Centrality hat viele Nachbarn und spielt daher eine wichtige Rolle für die lokale Konnektivität des Netzwerks. In ungerichteten Netzwerken entspricht die Degree Centrality der Anzahl aller benachbarten Knoten. (Freeman, 1979).
\subsubsection{Topologien}
Komplexe Netzwerke können auf unterschiedliche Weise strukturiert sein. Drei grundlegende Modelle haben sich bisher zur Beschreibung und Simulation realer Systeme etabliert: Erdős{–}Rényi{-}Netzwerke, Watts{–}Strogatz-Small-World-Netzwerke und Barabási{–}Albert{-}skalenfreie Netzwerke. Diese Topologien unterscheiden sich hinsichtlich Verbindungsstruktur, Gradverteilung, Clusterkoeffizient und typischer Pfadlängen und damit auch hinsichtlich ihrer Resilienz.
Das Erdős{–}Rényi-Netzwerk, ER-Modell, (Erdős & Rényi, 1959) erzeugt ein Netzwerk, indem zwischen allen möglichen Knotenpaaren jeweils mit einer festen Wahrscheinlichkeit p eine Kante gezogen wird. Diese zufällige Verbindungsstrategie hat mehrere Konsequenzen:
	Die Gradverteilung ist binomial- bzw. in großen Netzen näherungsweise Poisson-verteilt und die meisten Knoten haben ähnliche Gradwerte.
	Der Clusterkoeffizient ist vergleichsweise gering, da nur wenige Dreiecksbeziehungen entstehen.
	Die Pfadlänge ist typischerweise klein (logarithmisch mit Netzwerkgröße).
	Die Resilienz ist moderat robust gegenüber zufälligen Ausfällen, aber vulnerabel gegenüber gezielten Angriffen, da keine systematische Struktur existiert, die kritische Knoten schützt.
ER-Netzwerke bilden eine wichtige Baseline, aber reale Stromnetze weisen häufig strukturierte Muster auf, die durch dieses Modell nicht vollständig beschrieben werden können.
Das Small{-}World{-}Netzwerk{-}Modell (Watts & Strogatz, 1998) kombiniert Merkmale regulärer und zufälliger Netzwerke. Die Konstruktion erfolgt durch Ausgang eines regulären Rings, in dem Knoten mit ihren nächsten Nachbarn verbunden sind, und anschließend eine zufällige „Umverteilung“ einzelner Kanten passiert.
Charakteristische Eigenschaften:
	Eine hohe lokale Clusterbildung~-~also eine lokale Redundanz~-~wird geschaffen. Diese Knoten haben viele gemeinsame Nachbarn.
	Die durchschnittliche Pfadlänge ist eher kurz. Dadurch existiert eine effiziente Erreichbarkeit im gesamten Netzwerk.
	Die Resilienz ist durch hohe lokale Clusterung vergleichsweise hoch und damit das Netzwerk widerstandsfähiger gegenüber Ausfällen. Jedoch ist es anfällig für Störungen, wenn die wenigen „Shortcut“-Kanten ausfallen, die die Pfadlänge stark reduzieren.
Small-World-Strukturen treten in vielen Energie- und Informationsnetzen auf, da sie Effizienz mit lokaler Stabilität verbinden.
Von besonderer Bedeutung sind skalenfreie Netzwerke nach Barabási und Albert, deren Knotengradverteilung einem Potenzgesetz folgt. Das BA-Modell (Albert, Jeong & Barabási, 2000) basiert auf „präferentieller Anbindung“ von Knoten. Es werden nacheinander Knoten hinzugefügt und sie werden bevorzugt mit bereits gut vernetzten Knoten verbunden („the rich get richer“). Dadurch entstehen sog. Hubs.
Eigenschaften:
	Die Gradverteilung folgt einem Potenzgesetz (scale-free). Dadurch entstehen wenige Hubs mit vielen Knoten mit geringem Grad.
	Der Clusterkoeffizient ist vergleichsweise gering, kann aber je nach Erweiterung variieren.
	Die Pfadlängen ist sehr kurz, da Hubs viele Wege verbinden (ultra-small world).
	Die Resilienz ist hoch, da das Netzwerk sehr robust gegen zufällige Ausfälle ist (da diese meist kleine Knoten betreffen). Es ist aber stark anfällig gegenüber gezielten Angriffen auf Hubs, da das Netzwerk leicht fragmentieren kann.
Viele reale Infrastrukturnetze, einschließlich Kommunikationssysteme, Internet-Backbones und Teile moderner Stromnetze, zeigen skalenfreie Eigenschaften. Daher ist das BA-Modell besonders relevant für Analysen der Verwundbarkeit kritischer Infrastrukturen.

\subsection{Robustheit, Resilienz, Fehlertoleranz}
In der Literatur zu komplexen Netzwerken werden die Begriffe Robustheit, Fehlertoleranz und Resilienz häufig gemeinsam verwendet, obwohl sie sich auf unterschiedliche Aspekte der Systemstabilität beziehen. In den folgenden Unterkapiteln wird eine Abgrenzung im Kontext von Stromnetzwerken vorgenommen, basierend auf der Übersichtsarbeit von Xiaoyu Qi und Gang Mei (2024). 
\subsubsection{Robustheit}
Robustheit kann als Teilaspekt von Resilienz betrachtet werden  und bezieht sich hauptsächlich darauf, wie widerstandsfähig ein Netzwerk in seiner Struktur selbst gegenüber Störungen ist. Im Bereich von Stromnetzwerken, die zu den technologischen Netzwerken zählen, wird Robustheit als die Resistenz gegenüber Ausfällen bezeichnet. Die Kernfrage der Robustheit besteht darin, wie gut die Netzwerkfunktion unmittelbar nach einer Störung erhalten bleibt. In diesem Sinne gilt ein Stromnetz als robust, der Ausfall von einzelnen Komponenten wie Umspannwerken oder Leitungen nicht zu einer großflächigen Fragmentierung führt und ein Großteil der Knoten weiterhin über zusammenhängende Pfade erreicht werden kann. 
Gemessen wird die Robustheit typischer Weise über die größte zusammenhängende Komponente (GCC). Hierbei werden nach und nach Knoten oder Kanten entweder zufällig oder gezielt aus dem Netzwerk entfernt und nach jedem Schritt die Größe des größten funktionierenden Netzteils gemessen: 
\[
\mathrm{R} = \frac{1}{N} \sum_{q=1}^{N} S(q)
\]
Wobei R die Gesamtrobustheit eines Netzwerks bezeichnet und S(q) die relative Größe der größten Komponente nach q entfernten Knoten darstellt. 

\subsubsection{Fehlertoleranz und Resilienz}
Die Fehlertoleranz bezeichnet die Fähigkeit eines Systems, trotz Fehler einzelner Komponenten weiterhin funktionsfähig zu bleiben, ohne die Fehler zwingend reparieren zu müssen. Während Stromnetze im Bereich der technologischen Netze angesiedelt sind, bezieht sich die Fehlertoleranz mehr auf Informationsnetzwerke, zu denen Kommunikations- und Internetnetzwerke zählen. Dabei wird stärker auf Hard- und Software Bezug genommen und die Fähigkeit beschrieben, Fehler strukturell oder algorithmisch zu kompensieren, z.B. durch Redundanzen oder Umleitung von Flüssen. Die Kernfrage der Fehlertoleranz ist, wie gut das System während eines auftretenden Fehlers weiterarbeiten kann. So gilt ein Stromnetz als fehlertolerant, wenn bei einem Leitungsausfall automatisch alternative Wege für den Stromfluss gewählt werden. 
Da für die Analyse der Fehlertoleranz eine Modellierung von Flussdynamiken und Steuerungslogiken erforderlich ist, wird diese in der vorliegenden Arbeit nicht explizit untersucht. Sie stellt jedoch einen wichtigen ergänzenden Aspekt zur strukturellen Robustheit in Smart Grids dar. 

Die Resilienz umfasst die Widerstandsfähigkeit eines Netzwerks im Störungsfall, die Fähigkeit zur Begrenzung der Auswirkungen und die Fähigkeit zur Wiederaufnahme eines akzeptablen Funktionsniveaus. Dementsprechend ist die Resilienz ein übergreifendes Konzept, welches die Robustheit, Fehlertoleranz, Anpassungs- und Erholungsfähigkeit beinhaltet. Ein Smart Grid ist beispielsweise dann resilient, wenn es nach einem extremen Wetterereignis schnell wieder auf ein normales Versorgungsniveau zurückkehren kann. 
Die Kernfrage der Resilienz ist, wie viel der ursprünglichen Leistung noch vorhanden ist. Die Bestimmung der Resilienz erfordert neben strukturellen Modellen auch dynamische Modelle zur Erholung, Reparatur und zur Systemsteuerung. Diese umfassende Betrachtung geht über den Fokus der vorliegenden Arbeit hinaus. 
Zusammenfassend lässt sich sagen, dass Robustheit, Fehlertoleranz und Resilienz verschiedene Perspektiven auf die Stabilität von Stromnetzen und Smart Grid Netzen bieten. Der Fokus dieser Arbeit liegt jedoch auf struktureller Robustheit, um zu analysieren, wie die Netzwerkstruktur sowie die Interdependenzen zwischen Kommunikations- und Stromnetz die Anfälligkeit gegenüber Angriffen beeinflussen. 
